\section{A proposed research program}
\label{sec:research-program}

The next work should target proof obligations exposed by the construction,
not merely increase the number of heuristic moves attempted.  We propose
the following program, arranged from the nearest implementable extension
to the missing global complexity theorem.  Each question includes a
concrete outcome by which progress can be judged.

\subsection{Immediate geometric and implementation questions}

\paragraph{1. Can a canonical raw obstruction guide geometric normalization?}
The residual stopping set identifies precisely which original and saturated
rows cannot disappear in the fixed model.  For the eight retained native
positive cuts with raw obstructions, locate a source-authenticated move
that changes one of these rows' actual cocore intersections.  Candidates
include boundary bigon removal, a different meridian disc system, and a
small certified change of the core handle structure.  The outcome should
be a before/after geometric certificate and a strictly smaller stated
measure, rather than a freely reduced relator that happens to contain
a singleton.  This is a finite, already supplied test set on which a
successful next theorem would have visible consequences.

\paragraph{2. How can geometric unit intersections be discovered efficiently?}
An algebraic exponent of one is insufficient.  One needs an embedded
attaching curve, an actual meridian disc, and a controlled isotopy or disc
slide proving one geometric intersection.  Lackenby's curve algorithms
compute geometric intersection and isotopy information in polynomial
time in surface triangulation size and logarithmic curve weights
\cite{lackenby-curves}.  The research task is to connect their exact input
and output structures to the source boundary charts here.  A separate
minimum-intersection calculation for each pair of curves would not by
itself certify one simultaneous complete meridian system.

\paragraph{3. Which three-handle cores admit compressed cancellation?}
Seven of the ten retained three-handle cases are actual unions of
handlebodies according to the explicit oracle.  Develop a source verifier
for a safe two/three-handle cancellation or a controlled retriangulation,
including how every bundle port is transported.  Start with these seven
cases and retain the three negative controls.  A sphere boundary or an
Euler calculation must not authorize deletion of a closed summand.
An eventual theorem should bound the number of new source records, not
only the change in Euler characteristic.

\paragraph{4. Can private cocores be addressed without being enumerated?}
The present theorem supplies their count and initial incidence type, but
not a name for every cocore in a genus-$2^B$ bundle.  A useful extension
would provide a parameterized family of cocores with a bounded number of
chart types, a binary address within each family, and independently
checkable intersection queries.  Its output contract should allow finite
queries rather than require a literal list of all handles.  The initial
target can be patterns using only polynomially many source ports; the
general case must allow the possibility of larger representation costs.

\paragraph{5. What boundary-pattern information survives the compression?}
An unmarked genus histogram omits cyclic order, endpoints, coorientation,
and the placement of graph vertices.  Design a richer certificate that
retains this data through the collar construction, connector contractions,
and geometric cancellations.  Matching the genus and a matrix of unsigned
intersections is not an adequate test.  A strong validation target is an
independent reconstruction of the induced pattern on small explicit cut
boundaries, followed by exact agreement of essential-pattern tests.

\paragraph{6. Can the finite compiler be made linear in its explicit input?}
The new peeling pass is linear in incidence work, but the delivered
relative component assembly and independent forest reconstruction are
deliberately straightforward polynomial routines.  Index each row, port,
and base once by its component, and replace repeated scans by one grouped
pass.  Preserve the independent checker architecture rather than importing
the producer's union--find routine.  Measure the complete verified task:
the current common-success timings show that optimizing a small finite
subroutine may have little effect when orbit discovery dominates.

\subsection{Compressed geometric transitions}

\paragraph{7. Can the incoming cut grammars compile actual hierarchy moves?}
The incoming tile and first-return results already evaluate important
compressed local queries\cite{incoming-disconnected,incoming-sparse,incoming-moments}.
The next theorem should start with authenticated source charts and end
with the exact charts of the modified manifold, including the new boundary
pattern.  Each grammar edge needs a geometric carrier; each gluing needs
orientation and endpoint checks.  A polynomial evaluation bound for an
already supplied grammar is only one part of this transition theorem.

\paragraph{8. When does a sequential hierarchy admit an ambient block atlas?}
Theorem~\ref{thm:seam-repair} assumes an atlas; it does not construct one
from arbitrary later surfaces.  Characterize the obstruction to matching
opposite boundary traces when those surfaces are reglued into an earlier
manifold.  An algorithm should either produce a properly embedded ambient
family with verified local links, or exhibit the unmatched seam or
self-intersection that prevents such a family.  This makes the extra
hypothesis testable rather than treating it as an informal picture.

\paragraph{9. Can many seam exchanges be batched with a real decrease bound?}
The sharpened repair estimate is
$\lfloor C_0/4\rfloor+I_0$, and $I_0$ can be exponentially large in a
binary slope description.  Seek a theorem that removes a whole certified
parallel family of intersection circles in one compressed operation, with
an amortized bound on the number of family changes.  Merely storing $I_0$
in binary or taking its logarithm does not change the number of unit
repairs.  The output should include the new surface carriers and a proof
that no remaining seam is reconnected incorrectly.

\paragraph{10. Can normalization respect the block potential?}
Ordinary repair leaves standard embedded surfaces that need not be normal
in the current handle structure.  Normalizing one while ignoring the
others can increase intersection complexity.  Prove either simultaneous
normalization with controlled $\sum c_P+\iota$, or a second bounded
quantity paying for every increase.  The individually nonzero variant
of the repair theorem allows many labels even in a knot exterior, but
does not guarantee that every later cut-open surface is nonseparating.
That sequential admissibility condition must be checked separately.

\paragraph{11. Can the component-selection step be performed from orbit data?}
The repair proof selects a compression component with nonzero relative
class, or one outside the span of the other classes.  To implement that
selection without expanding sheets, enrich the component oracle with
authenticated relative homology weights and orientation data.  Verify
that the weights really add under regluing and that the selected component
has the claimed boundary carrier.  A vector of abstract homology labels
without such carriers would not establish the embedded atlas hypotheses.

\subsection{The genuinely global questions}

\paragraph{12. Can useful surfaces be found within the intended candidate budget?}
A short certificate and a fast cut consumer do not bound discovery.
Establish a coverage theorem for a specific family of normal or compressed
surfaces: every eligible hierarchy state must contain a useful candidate
within a stated, efficiently searchable family.  The theorem must control
both the number of alternatives and their encoded sizes.  This is the
place to connect the earlier sparse-support and tile kernels to the
selection problem, rather than treating those kernels as an enumeration
bound.

\paragraph{13. Is there a logarithmic-depth structural decomposition?}
To obtain the February exponent by
Theorem~\ref{thm:conditional-composition}, one sufficient target is
$L=O(\log n)$ with polynomial candidate and state bounds.  The March
exponent permits $L=O((\log n)^2)$ under the same hypotheses.  Seek a
structural operation that decreases a suitable charged complexity by a
fixed proportion, or an amortized substitute.  A proof that a positive
integer decreases by one gives termination and a linear numerical bound,
which is not a logarithmic-depth theorem.

\paragraph{14. Can backward corrections be charged without resetting a full suffix?}
The ambient-block theorem demonstrates one sufficient local condition
for retaining other surface choices.  The next global result would have
to decompose every relevant correction into such blocks, control the
number and cost of block changes, and cover pattern-compressions and
vertical-boundary moves.  The goal is an honest replacement for the
reset factor $(g+1)^L$, with a potential maintained across all those
operations.  Counterexamples in which a later surface cannot be reglued
should be collected alongside successful block repairs.

\paragraph{15. Can all encoded-size bounds be composed from the diagram?}
Maintain a ledger for triangulation records, pattern records, normal
coordinate bits, private-family descriptions, and certificate trace size
after every operation.  Prove a bound in the original crossing count,
including transitions that simplify topology but enlarge its encoding.
Only after discovery, transport, and reset costs fit one such ledger can
the local results legitimately enter a global quasi-polynomial theorem.
The final test is a proof of the hypotheses of
Theorem~\ref{thm:conditional-composition}, or a different complete cost
argument, rather than an extrapolation of the benchmark curves.

\subsection{Proposed order of work}

The nearest next contribution is geometric normalization of the eight
retained raw-positive failures, paired with a source-preserving
two/three-handle rule for the seven retained three-handle positives.
These tasks have concrete source cases, exact oracles at small scale,
and a clear success criterion.  Boundary-pattern transport and private
cocore addressing should proceed alongside them, because an unmarked
terminal result alone cannot drive the intended hierarchy.  Ambient
block construction and batched repairs then address the reset bottleneck.
Candidate coverage and global size/depth control remain the final
structural obligations.
